\chapter{Challenges and Solutions}
\label{ch:challenges}

This chapter goes through the main problems we met, in the order we met them.
For each one, we describe what went wrong, what caused it, what we tried, and
what finally worked. Some attempts failed, and we include those too, because
they explain why the final code looks the way it does.

\section{Inaccurate Autonomous Movement}
\label{sec:ch-geometry}

\textbf{Problem.} Every autonomous drive and turn ended in the wrong place, and
long drives stopped well short.

\textbf{Cause.} There were three separate causes, which made the problem hard to
spot because they all produced similar-looking errors:

\begin{enumerate}
  \item The \code{DriveTrain} used the SDK's example measurements instead of our
    robot's (\sref{sec:original}), so every conversion was off by several
    percent.
  \item \code{BRAKE} kept the overshoot from each move, so errors added up over a
    routine (\sref{sec:hold}).
  \item A flat 5-second timeout cut long drives off at about half distance
    (\sref{sec:timeouts}).
\end{enumerate}

\textbf{Solution.} We replaced the hidden conversion with our own maths using
measured sizes, switched to \code{HOLD}, and made timeouts depend on the move
(commit \commit{eb032a0}). For the error that remained, mostly wheel scrub, we
added \code{TURN_SCALE} and the on-robot \code{turn_tuner()}.

\section{One File on the Robot}
\label{sec:ch-onefile}

\textbf{Problem.} After we split the program into files, it worked in the editor
but crashed on the robot with an \code{ImportError}.

\textbf{Cause.} The Brain only receives and runs one file (\sref{sec:singlefile}).

\textbf{Solution.} \code{build.py} joins the files back together before every
download (\sref{sec:build}). The cost is one extra step: forgetting to run
\code{build.py} means the robot runs an old version. This was, in fact, the first
thing the README tells a reader to check when ``nothing happens at all''.

\section{The Silent Vision Failure}
\label{sec:ch-silent}

This was the hardest problem to diagnose, because the sensor gave no error at
all. It simply never saw anything.

\textbf{Attempt 1: print to the console.} The first diagnostic, \code{vision_report()},
used \code{print()}. However, \code{print()} sends text to the VEXcode console on
the laptop, which is only visible while the robot is connected by cable with the
console open. On the field, the robot showed nothing, so we could not tell
whether the sensor was failing or the program was not running at all.

\textbf{Attempt 2: draw on the Brain screen} (\commit{b4624fd}). We rewrote the
report to draw on the Brain's own screen, with a tick counter, the
\code{installed()} result, the frame timestamp, and the object counts
(\sref{sec:diagnostics}). The screen showed the tick climbing, \code{installed}
as \code{True}, and the timestamp changing, but zero colours and zero codes with
the cone directly in front of the camera. So the sensor was connected and sending
frames, but was not matching anything.

\textbf{Cause 1: detection switched off} (\commit{91a2514}). The SDK turns colour
processing off by default, and it must be switched on with
\code{color_detection(True)}. We added this, but the counts stayed at zero.

\textbf{Cause 2: the colours never reached the sensor} (\commit{08c26bb}, three
days later). We had originally chosen \emph{not} to declare the colours in code,
because declaring them sends new values to the sensor and we did not want to
overwrite the colours tuned in the AI Vision Utility. However, a screenshot of
the Utility showed the message ``No AI Vision Sensor Connected''. The colours had
been tuned while the sensor was not connected, so they were only ever saved in
the Utility, never on the sensor. The sensor was matching against an empty list.

\textbf{Solution.} We declared the colours in \code{config.py}, using RGB values
sampled from the Utility's swatches, so they are sent to the sensor every time
the program starts (\lref{lst:colours}). This removed the dependency on the
Utility entirely. It also revealed a smaller mistake: the Utility listed
\code{bluecone} first, so blue is colour 1 and yellow is colour 2, the opposite of
our earlier guess.

\textbf{Lesson.} A failure with no error message needs a diagnostic that rules
out one cause at a time. Each line on the vision report answers one question
(``is the program running?'', ``is the sensor plugged in?'', ``is it sending?'',
``is it matching?''), and together they narrowed the problem down to the one
place that was wrong.

\section{Screen Flicker}
\label{sec:ch-flicker}

\textbf{Problem.} The turn tuner's screen flickered constantly.

\textbf{Cause.} The tuner redrew the screen on every loop, 20 times per second,
and each redraw started by clearing the screen. So the screen was blank for part
of every redraw, which the eye sees as flicker.

\textbf{Solution} (\commit{2e1cf54}). There were two changes. First, the screen
is now only redrawn when a button actually changes something. Second,
\code{brain.screen.render()} switches the screen to \emph{double buffering}: the
program draws onto a hidden copy, and \code{render()} shows the finished copy all
at once, so the blank stage is never visible.

Note: double buffering applies to the whole screen from then on. Any other code
that draws must then also call \code{render()}, or its output never appears. This
is safe in the tuner only because \code{turn_tuner()} never returns, so nothing
else draws afterwards. This is recorded in \code{AGENTS.md} as a pitfall.

\section{Driving to the Cone}
\label{sec:ch-cone}

On 22 and 23 October we tried to make the autonomous routine find the cone and
drive up to it. This took 42 commits over two sessions. We describe it in
stages, since each stage was a response to the failures of the one before.

\subsection{Stage 1: Spinning in a Circle}

The first version (\commit{3bf3746}) drove in a slow circle until the target
colour appeared. The search speed was difficult to choose:

\begin{itemize}
  \item At 30\% and 60\% power on the two sides, the robot turned too fast and the
    camera image blurred, so it spun straight past the cone.
  \item At 1\% and 2\%, it was too weak to overcome friction and barely moved.
  \item 5\% and 10\% was the compromise (\commit{cb78f60}).
\end{itemize}

What to search for also changed several times: first the yellow colour, then the
yellow--blue code, then yellow again, and finally the code (\commit{4e0d8ed}). The
yellow colour alone is easier to detect but less specific; the code is specific
but drops out more often. This trade-off was only solved later, by the
colour-pair fallback (\sref{sec:colourpair}).

\subsection{Stage 2: Steering While Driving}

Next, the robot was made to approach the cone, steering towards it while
driving, and stopping once the cone appeared a certain width in the image.

\begin{itemize}
  \item It stopped too early, sometimes immediately. A minimum driving time was
    added (\commit{f05e383}), and the stopping width was raised from 100 to 180,
    then 250 pixels (\commit{80b5dc2}, \commit{5bc2a42}).
  \item It drifted off-centre while steering. We changed it to turn first, until
    centred, and only then drive straight (\commit{dc63501}), and tightened the
    centring tolerance from 20 to 5 pixels (\commit{091819b}).
  \item It still wobbled, because the image shakes while the robot moves and the
    readings jump around. We tried driving in short pulses (\commit{3559957}) and
    requiring several centred frames in a row (\commit{7ea5cea}).
  \item We then tried proportional steering (\commit{5bffee1}), the same method
    later used by the cone assist (\sref{sec:assist}).
\end{itemize}

\subsection{Stage 3: Stop, Measure, Turn}

The key change came with \commit{2ea38be}. Instead of steering from readings
taken while moving, the robot now:

\begin{enumerate}
  \item stops and waits 250~ms to stop rocking;
  \item takes 7 frames and uses the \emph{median} \code{centerX};
  \item converts the error in pixels into an angle;
  \item turns by that angle using the encoders (\sref{sec:turn}).
\end{enumerate}

The \emph{median} is the middle value once the readings are sorted. We used it
instead of the mean (average) because one wild reading can drag the mean a long
way, but barely moves the median~\cite{tukey1977}. For example, the readings 162,
163, 163, 164, 240 have a mean of 178 but a median of 163.

To convert pixels into degrees, we used the camera's horizontal field of view.
The camera sees about 74° across 320 pixels, so each pixel is about
$74/320 \approx 0.23°$. A cone 20 pixels right of centre therefore needs a turn of
about $20 \times 0.23 = 4.6°$ to the right.

Testing showed that the robot consistently ended up aimed slightly to the right.
This was because the camera is not mounted exactly on the robot's centre line,
so ``straight ahead'' is not exactly pixel 160. We moved
\code{VISION_CENTER_X} to 164 (\commit{be39e7a}), a correction of about
$4 \times 0.23 \approx 0.9°$. This value is still used by the cone assist today.

\subsection{Stage 4: Making It Robust}

On 23 October, the routine became: wait, scan, approach in steps, then align. The
table below lists each failure and its fix.

\begin{table}[h]
  \centering
  \small
  \begin{tabularx}{\textwidth}{@{}XXl@{}}
    \toprule
    Failure seen in testing & Fix & Commit \\
    \midrule
    A continuous spin went past the cone & Turn in steps, stopping to look after each & \commit{0165a87} \\
    Cone not seen from far away & Minimum size 20 $\to$ 6 px; colour-pair fallback & \commit{6c3a1b7} \\
    Overshooting while closing in & Approach in shrinking steps, re-aiming between them & \commit{c0304ef} \\
    No way to tell if the sensor was alive at the start & Live sensor status during the start delay & \commit{08bd56b} \\
    Wobbling left and right when close & Merge the colour pair into one box; looser final tolerance & \commit{8471b29} \\
    Pushing into the cone & Stop when the cone stops getting wider & \commit{8471b29} \\
    Cone lost partway & Go back to scanning & \commit{86d01a7} \\
    Brief dropouts ended the approach & Drive up to 3 steps ``blind'' & \commit{2c93591} \\
    Sudden false readings close up & Ignore a width jump of more than $2.5\times$ & \commit{ce2e710} \\
    \bottomrule
  \end{tabularx}
  \caption{Failures and fixes on 23 October.}
  \label{tab:oct23}
\end{table}

The step size during the approach shrank as the cone grew:
\begin{equation}
  \text{step} = \max\!\left(40,\; 400 \times \left(1 - \frac{w}{240}\right)\right)~\text{mm},
  \label{eq:step}
\end{equation}
where $w$ is the cone's width in pixels and 240 is the width at which the robot
counts as having arrived.

\textbf{Worked example:} When the cone is far away and 60 pixels wide, the step
is $400 \times (1 - 60/240) = 300$~mm. When it is close and 200 pixels wide, the
step is $400 \times (1 - 200/240) \approx 67$~mm. So the robot covers distance
quickly at first, then slows down to avoid hitting the cone.

The ``stop when the cone stops getting wider'' rule worked because, once the
robot touches the cone, driving further does not bring it any closer, so its width
in the image stops increasing. This detected contact without a touch sensor.

\section{The Decision to Change Strategy}
\label{sec:ch-decision}

On 25 October, the whole cone-hunting routine was replaced by the fixed S--A--B
path (\commit{cf781d1}). The commit does not record a reason, but the history
suggests three:

\begin{itemize}
  \item \textbf{Each fix revealed a new failure.} Over 42 commits, the failures
    moved from one stage to the next rather than going away. This suggests the
    approach was limited by how reliable the sensor readings were, not by any one
    bug.
  \item \textbf{It was slow.} Every look took 250~ms of settling plus 7 frames of
    50~ms, about 0.6~s, before the robot could move again. A full scan in 15° steps
    could take over 15~s, a large part of an autonomous period.
  \item \textbf{It depended on conditions we cannot control,} such as lighting and
    what else is on the field. A competition routine must work the same way every
    time.
\end{itemize}

However, the work was not wasted. The parts that worked well, proportional
steering, the camera offset, and the colour-pair fallback, became the cone assist
in driver control. Here, a human decides when to use it and can cancel it
instantly, so an occasional failure costs very little.

\section{Known Pitfalls}
\label{sec:gotchas}

Some problems are recorded in \code{AGENTS.md} as warnings, so that anyone
changing the code can avoid them:

\begin{itemize}
  \item \textbf{Cartridge mismatch.} If a motor is declared with the green 18:1
    cartridge but actually has a blue 6:1 cartridge fitted, every move overshoots
    by exactly $18/6 = 3$ times. If drives and turns are both badly wrong by the
    same factor, the cartridge should be checked before anything else.
  \item \textbf{Reversed right motor.} The right motor is declared reversed, so the
    minus sign in \code{turn()} is what makes the sides oppose. Removing either
    one breaks turning.
  \item \textbf{\code{TURN_SCALE} is changed by \code{turn_tuner()}} using
    \code{global}. This only works because both are in the same file; moving
    either breaks it.
  \item \textbf{The Brain screen does not clear itself.} Printing a shorter line over
    a longer one leaves the end of the old line visible, so lines must be cleared
    first.
\end{itemize}

\section{Summary}

\begin{enumerate}
  \item Movement errors had three causes, fixed with measured maths, \code{HOLD},
    and proportional timeouts (\sref{sec:ch-geometry}).
  \item The vision sensor failed silently because detection was off and the
    colours had never been loaded; on-screen diagnostics found both
    (\sref{sec:ch-silent}).
  \item A vision-driven autonomous was built in four stages, but was replaced by a
    fixed path because it could not be made reliable enough
    (\sref{sec:ch-cone}, \sref{sec:ch-decision}).
\end{enumerate}

We now evaluate the program as it stands.
